\begin{tikzpicture}[font=\scriptsize,>=Stealth,line width=0.5pt]

% =====================================================================
% (a) OMA
% =====================================================================
\begin{scope}[shift={(0,0)}]
  \node[tit] at (0,3.95) {(a) OMA — tách tài nguyên};
  \foreach \x/\y/\lbl/\c in {0.5/0.7/UE$_3$/blue, 2.2/0.7/UE$_4$/violet,
                             0.5/1.7/UE$_1$/orange, 2.2/1.7/UE$_2$/teal}{
     \fill[\c!18] (\x,\y) rectangle ++(1.7,1.0);
     \draw[\c!60!black,line width=0.6pt] (\x,\y) rectangle ++(1.7,1.0);
     \node[font=\scriptsize] at (\x+0.85,\y+0.5) {\lbl};
  }
  \draw[gray!60,->] (0.5,0.5) -- (4.1,0.5) node[anchor=north,font=\tiny,gray!50!black,midway,yshift=-2pt]{thời gian};
  \draw[gray!60,->] (0.35,0.7) -- (0.35,2.9) node[anchor=east,font=\tiny,gray!50!black,midway,rotate=90,yshift=6pt]{tần số};
  \node[note] at (0,0.15)
    {Mỗi UE chiếm một khối tài nguyên riêng. Không có nhiễu liên người dùng,
     nhưng mỗi UE chỉ dùng được $1/K$ tài nguyên.};
\end{scope}

% =====================================================================
% (b) SDMA
% =====================================================================
\begin{scope}[shift={(8.6,0)}]
  \node[tit] at (0,3.95) {(b) SDMA — tách theo không gian};
  % BS nhieu anten
  \node[bx,minimum height=0.9cm,minimum width=0.5cm] (bs) at (0.75,1.7) {BS};
  \foreach \dy in {-0.25,0,0.25}{
     \draw[blue!70!black,line width=0.6pt] (0.98,1.7+\dy) -- ++(0.22,0);
     \fill[blue!65] (1.20,1.7+\dy-0.07) -- (1.20,1.7+\dy+0.07) -- (1.36,1.7+\dy) -- cycle;
  }
  % 4 bup song
  \foreach \y/\lbl/\c in {2.90/UE$_1$/orange, 2.20/UE$_2$/teal, 1.20/UE$_3$/blue, 0.50/UE$_4$/violet}{
     \fill[\c!22,opacity=0.75] (1.40,1.7) -- (4.15,\y+0.30) -- (4.15,\y-0.30) -- cycle;
     \node[ue,fill=\c!20,draw=\c!60!black] at (4.45,\y) {\lbl};
  }
  \node[note] at (0,0.15)
    {BS dùng véc-tơ tiền mã hóa để tạo búp sóng riêng cho từng UE.
     \textbf{Cần nhiều anten.} Hệ SISO của luận văn chỉ có một anten nên
     không tạo được phân tách không gian.};
\end{scope}

% =====================================================================
% (c) NOMA
% =====================================================================
\begin{scope}[shift={(0,-5.6)}]
  \node[tit] at (0,3.95) {(c) NOMA — tách theo miền công suất};
  % chong cong suat
  \foreach \y/\h/\lbl/\c in {0.75/0.95/$p_1$/orange, 1.70/0.70/$p_2$/teal,
                             2.40/0.48/$p_3$/blue, 2.88/0.32/$p_4$/violet}{
     \fill[\c!22] (0.45,\y) rectangle ++(0.75,\h);
     \draw[\c!60!black,line width=0.5pt] (0.45,\y) rectangle ++(0.75,\h);
     \node[font=\tiny] at (0.825,\y+\h/2) {\lbl};
  }
  \node[font=\tiny,anchor=north,align=center] at (0.825,0.70) {chồng\\ công suất};
  \draw[gray!60,->] (0.30,0.75) -- (0.30,3.25) node[anchor=east,font=\tiny,gray!50!black,midway,rotate=90,yshift=6pt]{công suất};

  % chuoi SIC nhieu tang: mot cot doc, khong cat nhau
  \node[sm,anchor=west,minimum width=1.6cm] (n1) at (2.05,3.05) {Giải mã $s_1$};
  \node[sm,anchor=west,minimum width=1.6cm,fill=gray!22] (n2) at (2.05,2.55) {SIC};
  \node[sm,anchor=west,minimum width=1.6cm] (n3) at (2.05,2.05) {Giải mã $s_2$};
  \node[sm,anchor=west,minimum width=1.6cm,fill=gray!22] (n4) at (2.05,1.55) {SIC};
  \node[sm,anchor=west,minimum width=1.6cm] (n5) at (2.05,1.05) {Giải mã $s_k$};
  \foreach \a/\b in {n1/n2, n2/n3, n3/n4, n4/n5}{ \draw[->] (\a.south) -- (\b.north); }
  \draw[decorate,decoration={brace,amplitude=4pt},gray!70] (3.80,3.20) -- (3.80,0.92);
  \node[font=\tiny,anchor=west,align=left,gray!40!black] at (3.98,2.06) {tối đa\\ $K-1$\\ tầng SIC};

  \node[note] at (0,0.18)
    {Các UE dùng chung tài nguyên, phân biệt bằng mức công suất.
     UE kênh tốt phải giải mã và loại lần lượt tín hiệu của các UE khác.
     Thứ tự giải mã phụ thuộc độ mạnh kênh.};
\end{scope}

% =====================================================================
% (d) RSMA
% =====================================================================
\begin{scope}[shift={(8.6,-5.6)}]
  \node[tit] at (0,3.95) {(d) RSMA — chung + riêng};
  % tach thong diep
  \node[font=\tiny,anchor=east] at (0.42,2.95) {$W_k$};
  \foreach \i/\x/\c in {1/0.50/orange, 2/1.42/teal, 3/2.34/blue, 4/3.26/violet}{
     \fill[magenta!25] (\x,2.80) rectangle ++(0.34,0.32);
     \draw[magenta!70!black,line width=0.4pt] (\x,2.80) rectangle ++(0.34,0.32);
     \fill[\c!25] (\x+0.34,2.80) rectangle ++(0.48,0.32);
     \draw[\c!60!black,line width=0.4pt] (\x+0.34,2.80) rectangle ++(0.48,0.32);
     \node[font=\tiny,anchor=north] at (\x+0.41,2.74) {\i};
  }
  \node[font=\tiny,anchor=west,magenta!70!black] at (4.20,3.08) {phần chung $W_{c,k}$};
  \node[font=\tiny,anchor=west,gray!40!black] at (4.20,2.80) {phần riêng $W_{p,k}$};

  % gop thanh luong chung + cac luong rieng
  \node[sm,fill=magenta!15,anchor=west] (sc) at (0.50,1.95) {gộp $\to$ luồng chung $s_c$};
  \node[sm,anchor=west] (sp) at (3.45,1.95) {$s_1,\dots,s_4$};
  \draw[->] (1.00,2.72) -- (1.00,2.22);
  \draw[->] (4.00,2.72) -- (4.00,2.22);

  % may thu cua UE k : dung 1 tang SIC
  \node[sm,anchor=west] (r1) at (0.50,1.05) {Giải mã $s_c$};
  \node[sm,fill=magenta!12,anchor=west] (r2) at (2.15,1.05) {SIC (1 tầng)};
  \node[sm,anchor=west] (r3) at (4.00,1.05) {Giải mã $s_k$};
  \draw[->] (r1.east) -- (r2.west);
  \draw[->] (r2.east) -- (r3.west);
  \draw[->] (1.20,1.75) -- (1.20,1.28);
  \draw[->] (4.30,1.75) -- (4.30,1.28);
  \node[font=\tiny,anchor=north west,align=left] at (0.50,0.90)
       {$R_k=\underbrace{\eta_k R_c}_{\text{phần chung}}+\underbrace{R_{p,k}}_{\text{phần riêng}}$};
  \node[note] at (0,0.15)
    {Thông điệp mỗi UE được tách thành phần chung và phần riêng.
     \textbf{Mọi UE giải luồng chung trước, rồi dùng đúng một tầng SIC}
     để loại $s_c$, sau đó giải luồng riêng của mình.};
\end{scope}

% =====================================================================
% Dai quan he bao ham
% =====================================================================
\begin{scope}[shift={(0,-8.05)}]
  \draw[gray!45,line width=0.9pt,->] (1.6,0) -- (12.2,0);
  \foreach \x in {2.0,11.8}{ \draw[gray!60] (\x,-0.12) -- (\x,0.12); }
  \node[font=\tiny,anchor=north] at (2.0,-0.18) {$\rho_c\to 0$};
  \node[font=\tiny,anchor=north] at (11.8,-0.18) {$\rho_c\to 1$};
  \node[bx,fill=teal!12,anchor=north] at (2.0,-0.55) {SDMA\\ \tiny chỉ luồng riêng};
  \node[bx,fill=violet!12,anchor=north] at (11.8,-0.55) {NOMA\\ \tiny gần như chỉ luồng chung};
  \fill[magenta!14] (5.0,-0.30) rectangle (9.0,0.30);
  \draw[magenta!60!black,line width=0.6pt] (5.0,-0.30) rectangle (9.0,0.30);
  \node[font=\scriptsize,magenta!60!black] at (7.0,0) {RSMA};
  \node[font=\tiny,anchor=north,align=center] at (7.0,-0.55)
       {RSMA bao hàm SDMA và NOMA như hai trường hợp đặc biệt;\\
        luận văn giới hạn miền hành động ở $\rho_c\in[0{,}70;\,0{,}99]$};
  \node[font=\tiny,anchor=east] at (1.5,0) {$\rho_c=p_c/P_{\max}$:};
\end{scope}

\end{tikzpicture}
